\documentclass[10pt]{exam} % Usa la clase 'exam'

% Configuración de la cabecera
\usepackage[utf8]{inputenc}
\usepackage[T1]{fontenc}
\usepackage{amsmath} % Por si necesitas matemáticas

\title{Examen Corto: Conceptos Fundamentales}
\author{Clase de Prueba}
\date{\today}

\begin{document}

\maketitle % Muestra el título, autor y fecha

% Puntuación total (opcional, pero útil)
\addpoints

\begin{questions} % Inicia el entorno de preguntas

% --- Pregunta 1: Opción Múltiple (2 Puntos) ---
\question[2] ¿Cuál es el comando utilizado en LaTeX para iniciar un nuevo párrafo sin sangría?
\begin{choices}
    \choice \verb+\newline+
    \choice \verb+\par+
    \choice \verb+\noindent+ % Respuesta correcta
    \choice \verb+\vspace+
\end{choices}

% --- Pregunta 2: Opción Múltiple (3 Puntos) ---
\question[3] ¿Qué estructura básica define un documento en el sistema de gestión de paquetes (Vim-Plug, LazyVim, etc.)?
\begin{choices}
    \choice init.lua
    \choice \verb+\documentclass{...}+ % Respuesta correcta
    \choice \verb+\begin{document}+
    \choice .config/nvim
\end{choices}

% --- Pregunta 3: Pregunta Abierta (5 Puntos) ---
\question[5] Utilizando notación matemática simple, explique brevemente la diferencia entre un editor de texto (como Neovim) y un IDE (como VS Code). Limite su respuesta a un párrafo.

\vspace{2in} % Deja 2 pulgadas de espacio para escribir la respuesta

\end{questions}

% Muestra la puntuación total al final
\end{document}